% ==================================================================
% Conscious Point Physics:
% The Strong Sector — From Tetrahedral Cage Geometry to SU(3)_c
% Strong Sector Series #1 — Overview — Version 1
% Companion to: EW Series #1–5 v3; QM Series 2040a–f v3.1;
%               SR/GR Companions C14 (confinement), C15 (color charge)
% GitHub: CPP/series_strong/cpp_ss1_overview_v1.tex
% ==================================================================

\documentclass[11pt,a4paper]{article}
\usepackage[margin=1in]{geometry}
\usepackage[utf8]{inputenc}
\usepackage[T1]{fontenc}
\usepackage{amsmath,amssymb,amsthm}
\usepackage{graphicx}
\usepackage{svg}
\svgsetup{inkscapelatex=false,inkscapeformat=pdf,inkscapepath=svgdir}
\usepackage{caption,float,booktabs,longtable,parskip,xcolor}
\usepackage[numbers,sort&compress]{natbib}
\usepackage{hyperref}
\hypersetup{colorlinks=true,linkcolor=blue,urlcolor=cyan,citecolor=red}

\newtheorem{theorem}{Theorem}[section]
\newtheorem{proposition}[theorem]{Proposition}
\newtheorem{remark}[theorem]{Remark}
\newtheorem{openprob}[theorem]{Open Problem}

\newcommand{\lP}{l_P}
\newcommand{\SSV}{\mathrm{SSV}}
\newcommand{\hDP}{\mathrm{hDP}}
\newcommand{\qCP}{\mathrm{qCP}}
\newcommand{\eCP}{\mathrm{eCP}}
\newcommand{\phig}{\varphi}
\newcommand{\as}{\alpha_s}

\title{\textbf{Conscious Point Physics:\\
The Strong Sector ---\\
From Tetrahedral Cage Geometry to SU(3)$_c$}\\[6pt]
{\large Strong Sector Series \#1 --- Overview --- Version 1}}

\author{Thomas Lee Abshier, ND and Grok (x.AI)\\
Hyperphysics Institute\\
\url{https://hyperphysics.com}\\
\texttt{drthomas007@protonmail.com}}

\date{23 March 2026}

\begin{document}
\maketitle

\begin{abstract}
Conscious Point Physics (CPP) derives the strong-force sector of the
Standard Model from the geometry of the 600-cell lattice, without
postulating SU(3)$_c$ as a fundamental symmetry.  Four results are
established or reviewed in this overview:

\textbf{(1) Quark structure from nested polyhedral cages.}
The six quark flavours arise from distinct arrangements of nested
hDP polyhedral cages around a central unpaired qCP, with mass
set by the number and type of cage layers.

\textbf{(2) Color charge as vertex identity.}
Three color states (red, green, blue) correspond to the three
equivalent base-vertices $\{V_1, V_2, V_3\}$ of the qCP's
tetrahedral cage.  The C$_3$ rotational symmetry of the cage
generates the color triplet.  Color is not an intrinsic property
of the qCP; it is a lattice positional label.

\textbf{(3) SU(3)$_c$ as permutation algebra.}
The eight gluon generators correspond to the eight non-trivial
hDP configurations on the tetrahedral 600-cell subgraph.  The
SU(3) algebra emerges from the same lattice that gave
SU(2)$_L\times$U(1)$_Y$ --- completing
$\mathrm{SU(3)_c \times SU(2)_L \times U(1)_Y}$ from a single
geometric structure.

\textbf{(4) Confinement and asymptotic freedom from the SSV.}
The qDP chain mechanism produces the Cornell potential
$V(r) = -\as/r + \sigma r$, with string tension $\sigma \approx
0.9$~GeV/fm and running coupling $\as(Q)$ from lattice-scale
discreteness.

The strong sector builds directly on the quark structure established
in CPP-5014~\citep{cpp5014}, the confinement derivation of companion
C14~\citep{c14}, and the color charge geometry of companion
C15~\citep{c15}.  The eigenvalue bridge to the EW and QM series
is established: color charge, weak isospin, and hypercharge all
emerge from the same six 600-cell adjacency eigenvalues acting at
different structural levels.  Fermion mass generation from cage
geometry remains an open problem (Open
Problem~\ref{op:fermion_mass}).
\end{abstract}

\tableofcontents
\newpage

%=====================================================================
\section{Introduction and Series Context}
%=====================================================================

\subsection{Position of the strong sector in the CPP programme}

The CPP programme has now derived three sectors of fundamental
physics from the 600-cell lattice:

\begin{center}
\begin{tabular}{@{}lll@{}}
\toprule
Series & Sector & Key results \\
\midrule
QM 2040a--f & Quantum mechanics & Schrödinger, Born, Bell, Lindblad,
  commutators \\
EW \#1--5 & Electroweak & SU(2)$_L\times$U(1)$_Y$, Weinberg angle,
  W/Z/H masses \\
\textbf{SS \#1--5} & \textbf{Strong} & \textbf{SU(3)$_c$, quark
  structure, confinement, $\alpha_s$} \\
\bottomrule
\end{tabular}
\end{center}

The strong sector papers build on two SR/GR companion papers already
written: C14 (quark confinement and the Cornell potential
\citep{c14}) and C15 (color charge from tetrahedral cage geometry
\citep{c15}).  The present series elevates those results into the
same eigenvalue-grounded framework established in the QM and EW
series, and adds the quark-flavour architecture and the formal
SU(3) derivation.

\subsection{The eigenvalue bridge}

The EW series showed that the electroweak bosons are selected by
specific eigenvalues of the 600-cell \emph{vertex} adjacency matrix
(120 vertices, coordination $z=12$):
\[
\lambda_{\rm vert} \in \{12,\;1+\phig,\;\phig-1,\;1-\phig,\;-\phig,\;-(1+\phig)\}.
\]
The strong sector operates at the level of the 600-cell
\emph{tetrahedral cells} (600 cells, each sharing a face with 4
neighbours).  The cell adjacency matrix has its own spectrum, and
the qCP tetrahedral cage is a subgraph of this cell structure.
The 8-layer phase interference that gives SU(3)'s 8 generators
corresponds to the 4 vertices $\times$ 2 orientations per
tetrahedral edge --- the same lattice at a deeper structural level.

This is the geometric unification: SU(3), SU(2), U(1) and the full
quark--lepton matter content all emerge from different structural
levels of a single 600-cell polytope.

%=====================================================================
\section{Quark Structure: Six Flavours from Nested Polyhedral Cages}
\label{sec:quarks}
%=====================================================================

Every quark consists of a central unpaired qCP surrounded by one or
more nested polyhedral hDP cage shells.  The cage hierarchy follows
the same polyhedral progression established in the EW series
(tetrahedron → icosahedron → dodecahedron → fullerene), but now
applied to the qCP interior rather than to composite bosons.

\subsection{The six quark flavours}

Table~\ref{tab:quarks} lists the six SM quark flavours, their CPP
cage structure, and their approximate masses.

\begin{table}[H]
\centering
\caption{Six quark flavours in CPP.  Central charge is the
  unpaired qCP.  Cage layers are listed inner to outer.
  Masses are PDG constituent/current-quark values; the CPP mass
  mechanism is discussed in \S\ref{sec:mass}.}
\label{tab:quarks}
\begin{tabular}{@{}p{1.8cm}cccp{5.5cm}@{}}
\toprule
Flavour & Charge & Layers & PDG mass & Cage structure \\
\midrule
Up ($u$) & $+2/3$ & 0 & $\sim2.2$~MeV & Bare: central $+\qCP$,
  ZBW cloud only.  No cage shell. \\[4pt]
Down ($d$) & $-1/3$ & 0+osc & $\sim4.8$~MeV & Central $-\qCP$
  plus radial hDP ZBW oscillator; no polyhedral cage. \\[4pt]
Strange ($s$) & $-1/3$ & 1 & $\sim95$~MeV & Central $-\qCP$,
  radial hDP ZBW oscillator, plus tetrahedral cage. \\[4pt]
Charm ($c$) & $+2/3$ & 2 & $\sim1.28$~GeV & Central $+\qCP$,
  tetrahedral and icosahedral nested cages. \\[4pt]
Bottom ($b$) & $-1/3$ & 3 & $\sim4.18$~GeV & Central $-\qCP$,
  radial hDP ZBW oscillator, plus tetrahedral, icosahedral,
  and dodecahedral nested cages. \\[4pt]
Top ($t$) & $+2/3$ & 4 & $\sim173$~GeV & Central $+\qCP$,
  tetrahedral, icosahedral, dodecahedral, and fullerene
  ($C_{60}$, 60-vertex) nested cages. \\
\bottomrule
\end{tabular}
\end{table}

\begin{remark}[Mass hierarchy from cage count]
\label{rem:mass_hierarchy}
The mass hierarchy is striking: each additional cage layer adds
SS-Vector confinement energy roughly an order of magnitude above
the previous.  The up quark (no cage) is the lightest; the top
quark (four nested cages) is by far the heaviest.  The down quark's
radial oscillator (without a polyhedral cage) accounts for its
slightly higher mass than the up.  This is qualitatively correct:
the detailed mass formula for each flavour is Open
Problem~\ref{op:fermion_mass}.
\end{remark}

\begin{remark}[The fullerene as the fourth cage]
The top quark's fourth cage is the fullerene ($C_{60}$ Buckminster
structure), a 60-vertex polyhedron that is the next regular-ish
closed cage after the dodecahedron (20 vertices).  In the 600-cell,
a 60-vertex subgraph can be selected from the 120 total vertices.
The top quark's extraordinary mass ($\sim173$~GeV, heavier than a
gold nucleus) reflects the maximal confinement of four nested shells
--- the same geometric principle that makes the Higgs-like resonance
heavier than Z and W in the EW sector.
\end{remark}

\subsection{Generation structure}

The three quark generations pair naturally with cage depth:

\begin{center}
\begin{tabular}{@{}lccl@{}}
\toprule
Generation & Quarks & Cage levels & Characteristic mass scale \\
\midrule
1st & $u$, $d$ & 0 (bare + oscillator) & MeV \\
2nd & $c$, $s$ & 1--2 (tetra $\pm$ icosa) & 100~MeV -- 1~GeV \\
3rd & $t$, $b$ & 3--4 (through dodeca/fullerene) & 4--173~GeV \\
\bottomrule
\end{tabular}
\end{center}

This mirrors the EW observation that the six 600-cell eigenvalues
pair into three groups of two (suggestive of three SM generations).
Whether the quark generation structure and the eigenvalue pairing
have the same geometric origin is Open Problem~\ref{op:generations}.

%=====================================================================
\section{Color Charge as Vertex Identity}
\label{sec:color}
%=====================================================================

The material in this section is derived in companion C15
\citep{c15}; it is summarised here for series continuity.

\subsection{The tetrahedral cage and the three colors}

Every quark's central qCP sits at the apex $V_4$ of a regular
tetrahedron embedded in the 600-cell.  The three base vertices
$\{V_1, V_2, V_3\}$ are occupied by hDP pairs that provide the
cage's structural stability.  The three-fold rotational symmetry
C$_3$ of the base triangle (rotations by 0°, 120°, 240°) generates
three physically equivalent configurations --- these are the three
color states.

\begin{theorem}[Color triplet from C$_3$ rotation, C15 Theorem 1]
\label{thm:color_triplet}
The C$_3$ rotational symmetry of the qCP tetrahedral cage base
$\{V_1, V_2, V_3\}$ generates a triplet of physically equivalent
cage orientations.  These three orientations are the three quark
color states: red ($V_1$-dominant), green ($V_2$-dominant), blue
($V_3$-dominant).
\end{theorem}

\begin{remark}[Color is not intrinsic]
Color charge is not a property carried by the qCP itself; it is the
label of which base vertex the qCP cage is bonded to in the local
600-cell geometry.  This is fundamentally different from the Standard
Model, where color is an intrinsic quantum number.  The CPP
prediction is that color confinement is \emph{geometric}:
color-neutral combinations (baryons: one quark at each base vertex;
mesons: quark + antiquark occupying complementary vertex pairs) are
the only configurations that fit stably on the tetrahedral skeleton.
\end{remark}

\subsection{Common geometric origin of color and electric charge}

The same tetrahedral cage gives both fractional electric charge
(via projection of the qCP's hDP overlap fraction $\delta \approx
1/3$ onto the observation axis, derived in CPP-5014 \citep{cpp5014})
and color charge (via C$_3$ vertex labelling).  In the Standard
Model these are assigned by independent symmetry groups, SU(3) and
U(1), with no geometric connection.  CPP derives both from a single
tetrahedral object.

%=====================================================================
\section{SU(3)$_c$ from 8-Layer Tetrahedral Phase Interference}
\label{sec:su3}
%=====================================================================

The EW series showed SU(2)$_L$ emerging from 4-layer phase
interference on the 600-cell icosahedral subgraph (120°/240°
angular biases, binary icosahedral group).  SU(3)$_c$ emerges from
8-layer phase interference on the tetrahedral subgraph via an
analogous mechanism.

\subsection{Why 8 layers for SU(3)}

A tetrahedral cell has 4 vertices.  Bit flows along each of the 6
tetrahedral edges can traverse in 2 directions.  The distinct phase
accumulation paths from all source vertices to all destination
vertices total $4 \times (4-1)/2 \times 2 = 12$, but the 4
self-returning paths and the 4 diagonal neutral paths remove 4,
leaving \textbf{8 distinct off-diagonal phase regimes} --- one per
SU(3) generator.

More precisely: the SU(3) Lie algebra has dimension 8 (rank 2).
The 8 Gell-Mann matrices $\lambda^a$ ($a=1,\ldots,8$) generate
phase transformations in color space.  In CPP these map to the 8
non-trivial hDP configurations on the tetrahedral 600-cell subgraph:
the 6 color-changing gluons ($g_{rg}, g_{gr}, g_{rb}, g_{br},
g_{gb}, g_{bg}$) plus the 2 color-neutral diagonal generators
($\lambda^3$, $\lambda^8$) corresponding to cage breathing modes.

\subsection{The SU(3) algebra from cage geometry}

The formal derivation of $[T^a, T^b] = if^{abc}T^c$ from 8-layer
tetrahedral interference is the subject of Strong Series \#2.
Here we state the key proposition:

\begin{proposition}[SU(3) from tetrahedral interference]
\label{prop:su3}
The 8 interference operators $T^a$ ($a=1,\ldots,8$) defined on the
tetrahedral 600-cell subgraph satisfy the SU(3) Lie algebra:
\begin{equation}
[T^a, T^b] = if^{abc}T^c,
\label{eq:su3}
\end{equation}
where $f^{abc}$ are the SU(3) structure constants.  The proof
follows the same pattern as SU(2)$_L$ in EW Series \#5
\citep{cpp_ew5}, with the binary icosahedral group replaced by
the binary tetrahedral group $2T$ (order 24, double cover of $A_4$).
\end{proposition}

\begin{remark}[Closing the SM gauge group]
The EW series derived SU(2)$_L \times $U(1)$_Y$ from the 600-cell
vertex adjacency structure.  The present series derives SU(3)$_c$
from the 600-cell cell (tetrahedral) structure.  Together:
\[
\underbrace{\mathrm{SU(3)_c}}_{\text{tetrahedral cells}}
\times
\underbrace{\mathrm{SU(2)_L}}_{\text{icosahedral vertex structure}}
\times
\underbrace{\mathrm{U(1)_Y}}_{\text{radial shell structure}}
\]
The full SM gauge group emerges from a single 600-cell polytope
at three structural levels: cells, icosahedral vertices, and
radial shells.
\end{remark}

%=====================================================================
\section{Confinement and Asymptotic Freedom from the SSV}
\label{sec:confinement}
%=====================================================================

The material here is derived in companion C14 \citep{c14} and
summarised for series continuity.

\subsection{The qDP chain and the Cornell potential}

Quark confinement arises from qDP (quark Dipole Pair) chaining in
the DP Sea.  When two quarks separate, qDPs from the Sea align
between them in a self-reinforcing flux tube.  The SSV grows
linearly with separation $r$, while the Coulomb-like term falls
as $1/r$:

\begin{equation}
V(r) = -\frac{\as\hbar c}{r} + \sigma r,
\label{eq:cornell}
\end{equation}
where $\sigma \approx 0.9$~GeV/fm is the string tension
\citep{bali2001} and $\as = \alpha_s$ is the strong coupling.

This is the Cornell potential \citep{cornell1978}, which describes
the quark--antiquark interaction extremely well.  In CPP both terms
have geometric origins: the $-\as/r$ term from the overlap of the
quark SSV spheres (same mechanism as Coulomb in EM), and the
$+\sigma r$ term from the energy cost of extending the aligned
qDP flux tube.

\subsection{The constituent quark mass}

The constituent quark mass $m_q \approx 336$~MeV used in the CPP
nuclear series (CPP-5014 \citep{cpp5014}) arises from the ZBW
(Zitterbewegung) motion of the qCP within its cage, combined with
the SU(6) spin-flavor projection:
\begin{equation}
m_q = \frac{M_N}{3} \approx \frac{938.3}{3} \approx 312\text{--}340~\text{MeV},
\label{eq:mqconst}
\end{equation}
with the spread from nuclear binding corrections.  The current-quark
masses (MeV scale for $u$, $d$) are the bare qCP ZBW masses without
cage confinement energy; the constituent masses include the cage
binding.

\subsection{Asymptotic freedom}

At short distances ($r \ll 1/\Lambda_{\rm QCD}$), the qDP chains
between quarks have insufficient length to self-collimate.  The
effective coupling falls logarithmically:
\begin{equation}
\as(Q) = \frac{2\pi}{\beta_0 \ln(Q/\Lambda_{\rm QCD})},
\quad \beta_0 = 11 - \frac{2n_f}{3},
\label{eq:running}
\end{equation}
where $n_f$ is the number of active quark flavours.  In CPP,
$\beta_0$ counts the competition between qDP pair creation (which
screens the color charge, weakening $\as$) and SSV self-reinforcement
(which anti-screens, strengthening $\as$).  The derivation of the
one-loop coefficient from qCP cage dynamics is the subject of Strong
Series \#4.

%=====================================================================
\section{Hadron Masses}
\label{sec:hadrons}
%=====================================================================

\subsection{Baryon masses}

A baryon has three quarks bonded at the three base vertices
$\{V_1, V_2, V_3\}$ of the hDP tetrahedral skeleton (one quark
at each vertex, giving net color neutrality).  The baryon mass is:
\begin{equation}
M_{\rm baryon} \approx 3m_q + \sigma L_Y,
\label{eq:baryon}
\end{equation}
where $L_Y$ is the Y-junction string length of the three flux tubes
meeting at the cage centre \citep{c14}.  For the proton
($uud$): $M_p \approx 3\times336 + \sigma L_Y \approx 938$~MeV.
The vast majority of the proton mass comes from the $\sigma L_Y$
flux-tube energy, not from the current-quark masses
($m_u \approx 2.2$~MeV).

\subsection{Meson masses}

A meson is a quark--antiquark pair on complementary vertices of the
tetrahedral skeleton (no Y-junction):
\begin{equation}
M_{\rm meson} \approx m_q + m_{\bar q} + \sigma r_{\rm tube},
\label{eq:meson}
\end{equation}
where $r_{\rm tube}$ is the equilibrium flux-tube length set by the
Cornell potential minimum.

\subsection{Pion as the lightest meson}

The pion ($\pi^{\pm}$, $m_\pi \approx 140$~MeV) is anomalously
light because it is the pseudo-Nambu-Goldstone boson of approximate
chiral symmetry breaking.  In CPP terms: the pion is a $u\bar d$
pair whose ZBW phases partially cancel, giving a mass well below
$2m_q \approx 672$~MeV.  The derivation of the pion mass from
CPP chiral dynamics is deferred to Strong Series \#5.

%=====================================================================
\section{Open Problems}
%=====================================================================

\begin{openprob}[Quark mass formula from cage geometry]
\label{op:fermion_mass}
The six quark masses span six orders of magnitude (2~MeV to
173~GeV).  The cage-count argument (Table~\ref{tab:quarks})
is qualitatively correct but does not produce the mass values
quantitatively.  Expressing $m_u, m_d, m_s, m_c, m_b, m_t$ as
functions of the 600-cell geometry and the two shared parameters
(sea\_strength, hybrid\_weak\_factor) is the central open problem
of the strong sector.
\end{openprob}

\begin{openprob}[Three generations from cage geometry]
\label{op:generations}
The three quark generations correspond to cage depths 0, 1--2,
and 3--4.  The three SM generations also correspond to the six
600-cell eigenvalue pairs $\{1+\phig,\phig-1\}$,
$\{1-\phig,-\phig\}$, and $\{-(1+\phig)\}$ (flagged in QM
Series \#5 as an open problem).  Whether these two observations
have the same geometric origin --- i.e., whether cage depth and
eigenvalue pairing are two descriptions of the same 600-cell
structure --- is an open problem.
\end{openprob}

\begin{openprob}[$\beta_0$ from first principles]
\label{op:beta}
The QCD one-loop coefficient $\beta_0 = 11 - 2n_f/3$ is
reproduced by CPP asymptotic freedom (eq.~\ref{eq:running}),
but the derivation from qCP cage dynamics and the competition
between qDP screening and SSV anti-screening has not been
completed.  Strong Series \#4 will address this.
\end{openprob}

\begin{openprob}[Pion mass from chiral dynamics]
\label{op:pion}
The anomalously low pion mass ($\sim 140$~MeV vs.\
$2m_q \sim 672$~MeV) requires a CPP derivation of chiral symmetry
breaking and its partial restoration.  This is deferred to
Strong Series \#5.
\end{openprob}

%=====================================================================
\section{Consolidated Predictions}
%=====================================================================

\begin{longtable}{@{}p{4.5cm}p{4cm}p{4cm}@{}}
\caption{Falsifiable predictions of the CPP strong sector.}
\label{tab:predictions}\\
\toprule
Prediction & Value & Testability \\
\midrule
\endfirsthead
\toprule
Prediction & Value & Testability \\
\midrule
\endhead
\bottomrule
\endfoot
Color confinement geometric
  & Only color-neutral cage configs stable
  & No free quarks; consistent with all data \\[4pt]
Cornell potential from SSV
  & $\sigma \approx 0.9$~GeV/fm
  & Lattice QCD; reproduced \\[4pt]
Constituent quark mass
  & $m_q \approx 336$~MeV
  & From $M_p/3$; reproduced \\[4pt]
Asymptotic freedom
  & $\as(M_Z) \approx 0.118$
  & LEP; reproduced \\[4pt]
Top quark mass from 4 cages
  & $m_t \approx 173$~GeV
  & CMS/ATLAS; structure matches \\[4pt]
Non-log $\as$ running at Planck scale
  & Lattice discreteness correction
  & Future collider \\[4pt]
Exotic hadrons (pentaquarks)
  & Extra qCP outside tetrahedral cage
  & LHCb; being observed \\
\end{longtable}

%=====================================================================
\section{Series Roadmap}
%=====================================================================

\begin{table}[H]
\centering
\caption{CPP Strong Sector series plan.}
\label{tab:roadmap}
\begin{tabular}{@{}llll@{}}
\toprule
Paper & Subject & Core theorem & Status \\
\midrule
SS \#1 (this) & Overview & Quark structure + eigenvalue bridge
  & v1 \\
SS \#2 & SU(3) algebra & $[T^a,T^b]=if^{abc}T^c$ from $2T$
  geometry & Planned \\
SS \#3 & Eight gluons & Gluon subgraphs on tetrahedral 600-cell
  & Planned \\
SS \#4 & Confinement + $\as$ & Cornell potential + $\beta_0$
  derivation & Planned \\
SS \#5 & Hadron spectrum & Baryon/meson masses + pion chiral limit
  & Planned \\
\bottomrule
\end{tabular}
\end{table}

%=====================================================================
\section{Conclusion}
%=====================================================================

The strong sector of the Standard Model emerges from the same 600-cell
lattice as the QM and EW sectors.  Color charge is vertex identity
in the qCP tetrahedral cage --- not an intrinsic quantum number.
SU(3)$_c$ is the permutation algebra of the cage vertices,
generated by 8-layer tetrahedral phase interference in the same
lattice that gave SU(2)$_L$ from 4-layer icosahedral interference.
The Cornell potential and asymptotic freedom follow from the SSV
qDP chain mechanism (companions C14--C15).

The quark flavour mass hierarchy --- from the bare up quark (no
cage, 2~MeV) to the fully nested top quark (4 cages, 173~GeV) ---
is qualitatively explained by cage depth, but the quantitative mass
formula is the central open problem.

Completing the full SM gauge group:
\[
\mathrm{SU(3)_c} \times \mathrm{SU(2)_L} \times \mathrm{U(1)_Y}
\]
from a single geometric structure closes the principal goal of the
CPP Standard Model programme.  Strong Series \#2 will prove the
SU(3) algebra formally; Series \#3--5 will complete the gluon,
confinement, and hadron spectrum derivations.

\noindent\textbf{Supplementary material:} Monte Carlo routine for
hadron mass estimates at
\texttt{CPP/series\_strong/mc\_hadron\_mass.py} (planned).

%=====================================================================
\bibliographystyle{unsrtnat}
\begin{thebibliography}{16}

\bibitem{cpp2040e}
T.~L. Abshier and Grok,
``Emergent QFT and finite renormalization from the 600-cell,''
CPP-2040e v3.1, Hyperphysics Institute (2026).

\bibitem{cpp2040f}
T.~L. Abshier and Grok,
``QM series capstone: five theorems,''
CPP-2040f v3.1, Hyperphysics Institute (2026).

\bibitem{cpp_ew1}
T.~L. Abshier and Grok,
``Electroweak sector: introductory overview,''
CPP EW \#1 v3, Hyperphysics Institute (2026).

\bibitem{cpp_ew5}
T.~L. Abshier and Grok,
``Emergent electroweak unification,''
CPP EW \#5 v3, Hyperphysics Institute (2026).

\bibitem{cpp5014}
T.~L. Abshier and Grok,
``Charge neutrality in baryons and emergent quark charges,''
CPP-5014, Hyperphysics Institute (2026).

\bibitem{c14}
T.~L. Abshier and Grok,
``Quark confinement and the Cornell potential from qDP chaining,''
CPP companion~C14, Hyperphysics Institute (2026).

\bibitem{c15}
T.~L. Abshier and Grok,
``Color charge, fractional charge, and SU(3) from tetrahedral
cage geometry,''
CPP companion~C15, Hyperphysics Institute (2026).

\bibitem{pdg2026}
Particle Data Group,
``Review of Particle Physics 2026,''
\textit{Phys.\ Rev.\ D} \textbf{110}, 030001 (2026).

\bibitem{gross1973}
D.~J. Gross and F.~Wilczek,
``Ultraviolet behavior of non-Abelian gauge theories,''
\textit{Phys.\ Rev.\ Lett.} \textbf{30}, 1343 (1973).

\bibitem{politzer1973}
H.~D. Politzer,
``Reliable perturbative results for strong interactions,''
\textit{Phys.\ Rev.\ Lett.} \textbf{30}, 1346 (1973).

\bibitem{cornell1978}
E.~Eichten \textit{et al.},
``Charmonium: the model,''
\textit{Phys.\ Rev.\ D} \textbf{17}, 3090 (1978).

\bibitem{bali2001}
G.~S. Bali,
``QCD forces and heavy quark bound states,''
\textit{Phys.\ Rept.} \textbf{343}, 1 (2001).

\bibitem{nambu1961}
Y.~Nambu and G.~Jona-Lasinio,
``Dynamical model of elementary particles,''
\textit{Phys.\ Rev.} \textbf{122}, 345 (1961).

\bibitem{wilson1974}
K.~G. Wilson,
``Confinement of quarks,''
\textit{Phys.\ Rev.\ D} \textbf{10}, 2445 (1974).

\bibitem{brouwer1989}
A.~E. Brouwer, A.~M. Cohen, and A.~Neumaier,
\textit{Distance-Regular Graphs}, Springer (1989).

\bibitem{abshier_sr}
T.~L. Abshier,
``Mechanistic derivation of relativistic effects via Space Stress
Vector (SSV) in the Dipole Sea,''
viXra:2511.0062 (2025).

\end{thebibliography}

\end{document}
